\section{Debugging dan Alat Pengembang (DevTools)}

Browser menyediakan \textbf{DevTools} (F12) dengan beberapa panel penting untuk debugging JavaScript. Panel \textbf{Console} untuk menjalankan kode dan melihat log/error. Panel \textbf{Sources} untuk melihat kode, memasang breakpoint, dan step-through execution. Panel \textbf{Network} untuk memonitor HTTP request. Panel \textbf{Elements} untuk inspeksi DOM \cite{mdnToolsTesting}.

\begin{table}[htbp]
\centering
\begin{tabular}{|P{3cm}|P{4.5cm}|P{4.5cm}|}
\hline
\textbf{Teknik} & \textbf{Cara} & \textbf{Kapan Digunakan} \\
\hline
\code{console.log()} & Print nilai ke Console & Inspeksi nilai cepat \\
\hline
\code{console.table()} & Tampilkan data tabular & Array/objek \\
\hline
\code{debugger;} & Pause di baris ini & Breakpoint dalam kode \\
\hline
Breakpoint (UI) & Klik nomor baris di Sources & Debug tanpa ubah kode \\
\hline
Watch expression & Tambah di Sources panel & Monitor variabel \\
\hline
Network tab & Filter XHR/Fetch & Debug API call \\
\hline
\end{tabular}
\caption{Teknik debugging JavaScript dengan DevTools}
\end{table}

\begin{konsep}
\textbf{Breakpoint > console.log().} Meskipun \code{console.log()} cepat dan mudah, breakpoint lebih efektif untuk debugging kompleks: Anda dapat melihat seluruh call stack, semua variabel di scope saat ini, dan menjalankan kode step-by-step. Conditional breakpoint (klik kanan pada breakpoint) hanya pause jika kondisi terpenuhi---sangat berguna di loop besar \cite{mdnToolsTesting}.
\end{konsep}

